\documentclass[12pt]{article}

% Idioma y codificación
\usepackage[spanish, es-tabla]{babel}       %es-tabla para que se titule "Tabla"
\usepackage[utf8]{inputenc}

% Márgenes
\usepackage[a4paper,top=3cm,bottom=2.5cm,left=3cm,right=3cm]{geometry}

% Comentarios de bloque
\usepackage{verbatim}

% Paquetes de links
\usepackage[hidelinks]{hyperref}    % Permite enlaces
\usepackage{url}                    % redirecciona a la web

% Más opciones para enumeraciones
\usepackage{enumitem}


%------------------------------------------------------------------------

%Paquetes de figuras
\usepackage{caption}
\usepackage{subcaption} % Figuras al lado de otras
\usepackage{float}      % Poner figuras en el sitio indicado H.


% Paquetes de imágenes
\usepackage{graphicx}       % Paquete para añadir imágenes
\usepackage{transparent}    % Para manejar la opacidad de las figuras

% Paquete para usar colores
\usepackage[dvipsnames]{xcolor}
\usepackage{pagecolor}      % Para cambiar el color de la página

% Habilita tamaños de fuente mayores
\usepackage{fix-cm}


%------------------------------------------------------------------------

% Paquetes de matemáticas
\usepackage{amsmath, amsfonts, amssymb, mathrsfs}
\usepackage{cancel}     % Simplificar tachando
\usepackage{polynom}    % Divisiones y Ruffini

\usepackage{pgfplots}   %Representar funciones
\pgfplotsset{compat=1.18}  % Versión 1.18

%Definición de teoremas, etc.
\usepackage{amsthm}
%\swapnumbers   % Intercambia la posición del texto y de la numeración

\theoremstyle{plain}
\newtheorem{teo}{Teorema}[chapter]     % Se resetea en cada chapter
\newtheorem{coro}{Corolario}[teo]           % Se resetea en cada teorema
\newtheorem{prop}[teo]{Proposición}         % Usa el mismo contador que teorema
\newtheorem{lema}[teo]{Lema}                % Usa el mismo contador que teorema

\theoremstyle{remark}
\newtheorem*{observacion}{Observación}

\theoremstyle{definition}
\newtheorem{definicion}{Definición} [chapter]     % Se resetea en cada chapter
\newtheorem*{notacion}{Notacion}
\newtheorem*{ejemplo}{Ejemplo}
\newtheorem*{ejercicio*}{Ejercicio}             % No numerado
\newtheorem{ejercicio}{Ejercicio} [section]     % Se resetea en cada section


% Modificar el formato de la numeración del teorema "ejercicio"
\renewcommand{\theejercicio}{%
  \ifnum\value{section}=0 % Si no se ha iniciado ninguna sección
    \arabic{ejercicio}% Solo mostrar el número de ejercicio
  \else
    \thesection.\arabic{ejercicio}% Mostrar número de sección y número de ejercicio
  \fi
}


% \renewcommand\qedsymbol{$\blacksquare$}         % Cambiar símbolo QED
%------------------------------------------------------------------------

% Paquetes para encabezados
\usepackage{fancyhdr}
\pagestyle{fancy}
\fancyhf{}

\newcommand{\helv}{ % Modificación tamaño de letra
\fontfamily{}\fontsize{12}{12}\selectfont}
\setlength{\headheight}{15pt} % Amplía el tamaño del índice


%\usepackage{lastpage}   % Referenciar última pag   \pageref{LastPage}
\fancyfoot[C]{\thepage}

%------------------------------------------------------------------------

\begin{comment} % Descomentar en book. Comentar en article

% Conseguir que no ponga "Capítulo 1". Sino solo "1."
\makeatletter
\renewcommand{\@makechapterhead}[1]{%
  \vspace*{50\p@}%
  {\parindent \z@ \raggedright \normalfont
    \ifnum \c@secnumdepth >\m@ne
      \huge\bfseries \thechapter.\hspace{1em}\ignorespaces
    \fi
    \interlinepenalty\@M
    \Huge \bfseries #1\par\nobreak
    \vskip 40\p@
  }}
\makeatother
\end{comment}

%------------------------------------------------------------------------
% Personalizar la portada
\usepackage{titling}


%------------------------------------------------------------------------

% Comandos definidos
\newcommand{\bb}[1]{\mathbb{#1}}
\newcommand{\cc}[1]{\mathcal{#1}}
\newcommand{\Hop}{\stackrel{L'H\hat{o}pital}{=}}

\newcommand{\red}[1]{{\color{red}#1}} % Para integrales, destacar los cambios.

% Método de integración
\newcommand{\MetInt}[2]{
    \left[\begin{array}{c}
        #1 \\ #2
    \end{array}\right]
}

%------------------------------------------------------------------------
\newcommand{\T}[0]{\cc{T}}

\begin{document}

    % 1. Foto de fondo
    % 2. Título
    % 3. Encabezado Izquierdo
    % 4. Color de fondo
    % 5. Coord x del titulo
    % 6. Coord y del titulo
    % 7. Fecha

    
    % 1. Foto de fondo
% 2. Título
% 3. Encabezado Izquierdo
% 4. Color de fondo
% 5. Coord x del titulo
% 6. Coord y del titulo
% 7. Fecha

\newcommand{\portada}[7]{

    \portadaBase{#1}{#2}{#3}{#4}{#5}{#6}{#7}
    \portadaBook{#1}{#2}{#3}{#4}{#5}{#6}{#7}
}

\newcommand{\portadaExamen}[7]{

    \portadaBase{#1}{#2}{#3}{#4}{#5}{#6}{#7}
    \portadaArticle{#1}{#2}{#3}{#4}{#5}{#6}{#7}
}




\newcommand{\portadaBase}[7]{

    % Tiene la portada principal y la licencia Creative Commons
    
    % 1. Foto de fondo
    % 2. Título
    % 3. Encabezado Izquierdo
    % 4. Color de fondo
    % 5. Coord x del titulo
    % 6. Coord y del titulo
    % 7. Fecha
    
    
    \thispagestyle{empty}               % Sin encabezado ni pie de página
    \newgeometry{margin=0cm}        % Márgenes nulos para la primera página
    
    
    % Encabezado
    \fancyhead[L]{\helv #3}
    \fancyhead[R]{\helv \nouppercase{\leftmark}}
    
    
    \pagecolor{#4}        % Color de fondo para la portada
    
    \begin{figure}[p]
        \centering
        \transparent{0.3}           % Opacidad del 30% para la imagen
        
        \includegraphics[width=\paperwidth, keepaspectratio]{assets/#1}
    
        \begin{tikzpicture}[remember picture, overlay]
            \node[anchor=north west, text=white, opacity=1, font=\fontsize{60}{90}\selectfont\bfseries\sffamily, align=left] at (#5, #6) {#2};
            
            \node[anchor=south east, text=white, opacity=1, font=\fontsize{12}{18}\selectfont\sffamily, align=right] at (9.7, 3) {\textbf{\href{https://losdeldgiim.github.io/}{Los Del DGIIM}}};
            
            \node[anchor=south east, text=white, opacity=1, font=\fontsize{12}{15}\selectfont\sffamily, align=right] at (9.7, 1.8) {Doble Grado en Ingeniería Informática y Matemáticas\\Universidad de Granada};
        \end{tikzpicture}
    \end{figure}
    
    
    \restoregeometry        % Restaurar márgenes normales para las páginas subsiguientes
    \pagecolor{white}       % Restaurar el color de página
    
    
    \newpage
    \thispagestyle{empty}               % Sin encabezado ni pie de página
    \begin{tikzpicture}[remember picture, overlay]
        \node[anchor=south west, inner sep=3cm] at (current page.south west) {
            \begin{minipage}{0.5\paperwidth}
                \href{https://creativecommons.org/licenses/by-nc-nd/4.0/}{
                    \includegraphics[height=2cm]{assets/Licencia.png}
                }\vspace{1cm}\\
                Esta obra está bajo una
                \href{https://creativecommons.org/licenses/by-nc-nd/4.0/}{
                    Licencia Creative Commons Atribución-NoComercial-SinDerivadas 4.0 Internacional (CC BY-NC-ND 4.0).
                }\\
    
                Eres libre de compartir y redistribuir el contenido de esta obra en cualquier medio o formato, siempre y cuando des el crédito adecuado a los autores originales y no persigas fines comerciales. 
            \end{minipage}
        };
    \end{tikzpicture}
    
    
    
    % 1. Foto de fondo
    % 2. Título
    % 3. Encabezado Izquierdo
    % 4. Color de fondo
    % 5. Coord x del titulo
    % 6. Coord y del titulo
    % 7. Fecha


}


\newcommand{\portadaBook}[7]{

    % 1. Foto de fondo
    % 2. Título
    % 3. Encabezado Izquierdo
    % 4. Color de fondo
    % 5. Coord x del titulo
    % 6. Coord y del titulo
    % 7. Fecha

    % Personaliza el formato del título
    \pretitle{\begin{center}\bfseries\fontsize{42}{56}\selectfont}
    \posttitle{\par\end{center}\vspace{2em}}
    
    % Personaliza el formato del autor
    \preauthor{\begin{center}\Large}
    \postauthor{\par\end{center}\vfill}
    
    % Personaliza el formato de la fecha
    \predate{\begin{center}\huge}
    \postdate{\par\end{center}\vspace{2em}}
    
    \title{#2}
    \author{\href{https://losdeldgiim.github.io/}{Los Del DGIIM}}
    \date{Granada, #7}
    \maketitle
    
    \tableofcontents
}




\newcommand{\portadaArticle}[7]{

    % 1. Foto de fondo
    % 2. Título
    % 3. Encabezado Izquierdo
    % 4. Color de fondo
    % 5. Coord x del titulo
    % 6. Coord y del titulo
    % 7. Fecha

    % Personaliza el formato del título
    \pretitle{\begin{center}\bfseries\fontsize{42}{56}\selectfont}
    \posttitle{\par\end{center}\vspace{2em}}
    
    % Personaliza el formato del autor
    \preauthor{\begin{center}\Large}
    \postauthor{\par\end{center}\vspace{3em}}
    
    % Personaliza el formato de la fecha
    \predate{\begin{center}\huge}
    \postdate{\par\end{center}\vspace{5em}}
    
    \title{#2}
    \author{\href{https://losdeldgiim.github.io/}{Los Del DGIIM}}
    \date{Granada, #7}
    \thispagestyle{empty}               % Sin encabezado ni pie de página
    \maketitle
    \vfill
}
    \portadaExamen{ffccA4.jpg}{Topología I\\Examen IV}{Topología I. Examen IV}{MidnightBlue}{-8}{28}{2023-2024}{Arturo Olivares Martos}

    \begin{description}
        \item[Asignatura] Topología I.
        \item[Curso Académico] 2022-23.
        \item[Grado] Doble Grado en Ingeniería Informática y Matemáticas\footnote{El examen lo pone el departamento.}.
        \item[Grupo] Único.
        \item[Profesor] Miguel Ortega Titos.
        \item[Descripción] Convocatoria Ordinaria.
        \item[Fecha] 14 de enero de 2023.
        \item[Duración] 3 horas.
        \item[Observaciones] Todos los apartados de un mismo ejercicio tienen la misma puntuación.
    
    \end{description}
    \newpage
    
    \begin{ejercicio}[4.5 puntos]
        Para todo $R\geq 0$, consideramos el conjunto $S_R$ dado por $S_R= \{x\in \bb{R}^2\mid \|x\|= R\}$, donde $\|\cdot\|$ es la norma euclídea en $\bb{R}^2$.
        Se considera la topología $\T$ en $\bb{R}^2$ generada por la base $\cc{B}= \{S_R\mid R\geq 0\}$.
        \begin{enumerate}
            \item Estudiar cuándo el conjunto $\{x_0\}$, con $x_0\in \bb{R}^2$ arbitrario,
            es cerrado en $(\bb{R}^2,\T)$.

            \begin{itemize}
                \item Veamos en primer lugar que $\{0\}\in C_{\T}$. Tenemos que:
                \begin{equation*}
                    \bb{R}^2 \setminus \{0\} = \bigcup_{R>0} S_R \in \T
                \end{equation*}
                donde afirmamos que es un abierto por ser unión de abiertos. Por tanto, $\{0\}$ es cerrado.

                \item Veamos ahora que si $x_0\neq 0$, entonces $\{x_0\}$ no es cerrado.
                Sea $R_0=\|x_0\|$, y sea $x\in S_{R_0}$, $x\neq x_0$.
                Supongamos entonces que $\bb{R}^+\setminus \{x_0\}$ es abierto. Entonces,
                como $x\neq x_0$, existe $R>0$ tal que $x\in S_R\subset \bb{R}^+\setminus \{x_0\}$.
                Por tanto, como $x\in S_{R_0}$, tenemos que $x\in S_{R_0}\subset \bb{R}^+\setminus \{x_0\}$,
                algo que es una contradicción, ya que $x_0\in S_{R_0}$ y $x_0\notin \bb{R}^+\setminus \{x_0\}$.
            \end{itemize}


            \item Demostrar que $(\bb{R}^2,\T)$ es 1AN pero no es 2AN.
            
            \begin{itemize}
                \item Veamos en primer lugar que $(\bb{R}^2,\T)$ es 1AN. Sea $x\in \bb{R}^2$, y veamos que existe una base de entornos $\beta_x$ numerable.
                Sea $R_0=\|x\|$, y como tenemos que $\{B\in \cc{B}\mid x\in B\}$ es una base de entornos,
                tenemos que:
                \begin{equation*}
                    \beta_x = \{S_R\mid x\in S_R\} = \{S_R\mid \|x\|=R\} = \{S_{R_0}\}
                \end{equation*}
                Por tanto, dicha base de entornos no solo es numerable sino finita, con un único elemento.

                \item Veamos que no es 2AN. Sabemos que:
                \begin{equation*}
                    B(0,1) = \bigcup \{S_R\mid \|(0, x)\|=R,~x\in [0,1[\}
                \end{equation*}
                
                Supongamos que sí es 2AN, y llegaremos a una contradicción.
                Como es 2AN, entonces podemos extraer una base numerable de $\cc{B}$,
                que denotaremos por $\cc{B}'=\{S_{R_n}\}_{n\in \bb{N}}$. Como $B(0,1)$ es abierto, entonces
                \begin{equation*}
                    B(0,1) = \bigcup \{S_R\mid \|(0, x)\|=R,~x\in A,~A \text{ numerable.}\}
                \end{equation*}
                Como $[0,1[$ no es numerable, entonces $A\subsetneq [0,1[$. Por tanto, como tenemos que $A\neq [0,1[$,
                existe $x_0\in [0,1[\setminus A$, por lo que $S_{\|x_0\|}\subset \bb{R}^2\setminus B(0,1)$, lo que
                es una contradicción, ya que $S_{\|x_0\|}\subset B(0,1)$ por tener $\|x_0\|<1$.
            \end{itemize}

            \item Calcular la clausura, el interior y la frontera en $(\bb{R}^2,\T)$ del conjunto $A$
            dado por $A=\{(a,b)\in \bb{R}^2\mid |b|\leq 1\}$.
            \begin{itemize}
                \item Veamos en primer lugar que $A^{\circ} = \ol{B}(0,1)$.
                \begin{description}
                    \item[$\subset)$] Si $x\in A^{\circ}$, entonces existe $B\in \T$ tal que $x\in B\subset A$.
                    Como $\cc{B}$ da una partición de $\bb{R}^2$, entonces $B=S_{R_0}$, con $R_0=\|x\|$. Por tanto,
                    $S_{R_0}\subset A$, y como $S_{R_0}$ es cerrado, tenemos que $S_{R_0}\subset A$. Como $(0, R_0)\in S_{R_0}$,
                    tenemos que $(0, R_0) \in A$, por lo que $R_0\leq 1$. Por tanto, como $\|x\|\leq 1$, $x\in \ol{B}(0,1)$.

                    \item[$\supset)$] Si $(x,y)\in \ol{B}(0,1)$, entonces $x^2 + y^2 \leq 1$. Por tanto, $y^2\leq 1$, por lo que $|y|\leq 1$.
                    Por tanto, $(x,y)\in A$.
                \end{description}
            \end{itemize}

            \item Probar que $A$ es conexo en $(\bb{R}^2,\T)$ si y solo si existe $R\geq 0$ tal que $A\subset S_R$. Determinar
            las componentes conexas de $(\bb{R}^2,\T)$.

            Calculemos en primer lugar las componentes conexas de $(\bb{R}^2,\T)$. Sabemos que $\cc{B}= \{S_R\mid R\geq 0\}$ es una partición
            mediante abiertos de $\bb{R}^2$. Veamos ahora que $\left(S_R, \T_{S_R}\right)$ es conexo para todo $R\geq 0$. Calculamos en
            primer lugar $\T_{S_R}$, sabiendo que $\T=\{\ol{B}(0,r)\mid r>0\}$:
            \begin{equation*}
                U \in \T_{S_R} \iff \exists r'\in \bb{R}^+ \mid U = \ol{B}(0, r') \cap S_R = \ol{B}(0, r') \cap S(0,r)
            \end{equation*}
            Si $r'\geq r$, tenemos que $\ol{B}(0, r') \cap S(0,r) = S(0,r)$; mientras que si
            $r'<r$, entonces $\ol{B}(0, r') \cap S(0,r) = \emptyset$. Por tanto, $\T_{S_R} = \{\emptyset, S(0,r)\}$, que es la
            topología trivial. Por tanto, $\left(S_R, \T_{S_R}\right)$ es conexo,
            por lo que $\cc{B}$ es una partición mediante abiertos conexos de $\bb{R}^2$. Por tanto, las componentes conexas de
            $(\bb{R}^2,\T)$ son $\cc{B}$.\\

            Caracterizamos ahora los conexos de $(\bb{R}^2,\T)$ mediante doble implicación:
            \begin{description}
                \item[$\Longrightarrow)$] Sea $A$ conexo, y sea $x\in A$. Entonces, como $x\in A$ y $A$ conexo, $A\subset C(x)$.
                Por tanto, $A\subset S_R$ para algún $R\geq 0$.

                \item[$\Longleftarrow)$] Sea $R\geq 0$ tal que $A\subset S_R$. Veamos que $A$ es conexo. Como $A\subset S_R$,
                entonces $\T_{\big| A} = {\T_{S_R}}_{\big |A}$. Por tanto, ver que $\left(A, \T_A\right)$ es conexo  
                equivale a ver que $\left(A, {\T_{S_R}}_{\big |A}\right)$ es conexo. Como $\T_{S_R} = \{\emptyset, S(0,r)\}$, tenemos que es la topología
                trivial. Por tanto, como todo subconjunto de un espacio topológico trivial es conexo, tenemos que $\left(A, {\T_{S_R}}_{\big |A}\right)$
                es conexo, por lo que $A$ es conexo.
            \end{description}

            \item Demostrar que $A$ es compacto en $(\bb{R}^2,\T)$ si y solo si existe $J\subset [0,+\infty[$ finito tal que
            $A\subset \bigcup\limits_{R\in J}S_R$.

            \begin{description}
                \item[$\Longrightarrow)$] Sea $A$ compacto. Entonces, como $\cc{B}$ es un recubrimiento de $A$ por abiertos de $(\bb{R}^2,\T)$,
                tenemos que $\exists J\subset \bb{R}^+$ finito tal que $A\subset \bigcup\limits_{R\in J}S_R$.

                \item[$\Longleftarrow)$] Sea $J\subset \bb{R}^+$ finito tal que $A\subset \bigcup\limits_{R\in J}S_R$. Sabemos que $\left(S_R, \T_{S_R}\right)$ es la
                topología trivial. Por tanto, como todo subconjunto de un espacio topológico trivial es compacto, tenemos que $A\cap S_R$ es compacto. Veamos por tanto que
                $A$ es unión finita de compactos:
                \begin{equation*}
                    A = A \cap \left(\bigcup\limits_{R\in J}S_R\right) = \bigcup\limits_{R\in J} \left(A\cap S_R\right)
                \end{equation*}
                como $J$ es finito, tenemos que $A$ es unión finita de compactos, por lo que $A$ es compacto.
            \end{description}
        \end{enumerate}
    \end{ejercicio}


    \begin{ejercicio}[2.5 puntos] Enunciar y demostrar el teorema de Tichonov. Si se hace uso del lema del tubo,
        entonces éste debe ser enunciado y demostrado previamente.        
    \end{ejercicio}


    \begin{ejercicio}[3 puntos] Estudiar de forma razonada las siguientes cuestiones:
        \begin{enumerate}
            \item Decidir si los siguientes subespacios de $(\bb{R}^2,\T_u)$ son homeomorfos entre sí dos a dos o no:
            \begin{enumerate}
                \item $A_1=\{x\in \bb{R}^2\mid 1<\|x\|<4\}$,
                \item $A_2=\{x\in \bb{R}^2\mid 1\leq \|x\|\leq 4\}$,
                \item $A_3=\bb{R}^2\setminus \{0\}$.
            \end{enumerate}

            Veamos en primer lugar que $A_2$ es compacto. $$A_2=\{x\in \bb{R}^2\mid 1\leq \|x\|\}~\cap~\{x\in \bb{R}^2\mid \|x\|\geq 1\}$$
            Como $\|\cdot\|$ es continua con la topología usual, tenemos que $A_2$ es intersección de dos cerrados, luego es cerrado.
            Además, claramente $A_2\subset B(0, 5)$, luego $A_2$ es acotado. Por tanto, es compacto.

            Veamos ahora que $A_1$ no es compacto. Para ello, veremos que no es cerrado. Tenemos que:
            \begin{equation*}
                \bb{R}^2\setminus A_1 = \ol{B}(0,1) \cup \{x\in \bb{R}^2\mid \|x\|\geq 4\} 
            \end{equation*}

            Veamos que dicho conjunto no es abierto. Tomado $p=(0,1)\in \bb{R}^2\setminus A_1$,
            tenemos que no existe $U\in \T_u$, tal que $(0,1)\in U\cap \ol{B}(0,1)$.
            De forma directa se ve que $\nexists U\in \T$ tal que $(0,1)\in U\subset \bb{R}^2\setminus A_1$. Por tanto, $\bb{R}^2\setminus A_1$ no es entorno
            de $(0,1)$, por lo que $\bb{R}^2\setminus A_1$ no es abierto, y por tanto $A_1$ no es cerrado. Por tanto, tampoco es compacto.

            Además, como $A_3$ no está acotado, tenemos que no es compacto.\\

            Por tanto, tenemos que $A_2$ no es homeomorfo ni a $A_1$ ni a $A_3$. Veamos ahora que $A_1\cong A_3$.
            Sabemos que $]1,4[\cong ]0,\infty[$, por lo que sea $f$ el homeomorfismo entre ellos, y notamos $g=f^{-1}$. Definimos $\ol{f}:A_1\to A_3$ dado por:
            \begin{equation*}
                \ol{f}(x) = \frac{x}{\|x\|}f\left(\|x\|\right)
            \end{equation*}
            Sea $\ol{g}:A_3\to A_1$ dado por:
            \begin{equation*}
                \ol{g}(x) = \frac{x}{\|x\|} g\left(\|x\|\right)
            \end{equation*}
            Tenemos que ambas son continuas, ya que es el producto, cociente y composicicón de funciones continuas. Veamos que:
            \begin{align*}
                \ol{f}(\ol{g}(x)) &= \frac{\ol{g}(x)}{\|\ol{g}(x)\|}
                \cdot f\left(\|\ol{g}(x)\|\right)
                = \frac{\frac{x}{\|x\|} g\left(\|x\|\right)}{\left\|\frac{x}{\|x\|} g\left(\|x\|\right)\right\|}
                \cdot f\left(\left\|\frac{x}{\|x\|} g\left(\|x\|\right)\right\|\right) =\\
                &= \frac{\frac{x}{\|x\|} g\left(\|x\|\right)}{\frac{\|x\|}{\|x\|} \left\|g\left(\|x\|\right)\right\|}
                \cdot f\left(\left\|\frac{x}{\|x\|}\right\|g\left(\|x\|\right)\right) =
                \frac{x}{\|x\|}f(g(\|x\|)) \AstIg \frac{x}{\|x\|}\cdot \|x\| = x
            \end{align*}
            De igual forma, se demuestra que $\ol{g}(\ol{f}(x))=x~\forall x\in A_1$. De esta forma, se ve que $\ol{f},~\ol{g}$ son biyectivas,
            con $\ol{g}=\ol{f}^{-1}$. Por tanto, $\ol{f}$ es un homeomorfismo, por lo que $A_1\cong A_3$. 

            \item Sean $(X_1,\T_1), (X_2,\T_2)$ espacios topológicos. Para $i=1,2$, sea $\cc{R}_i$ una relación de equivalencia
            en $X_i$ tal que la proyección $p_i:(X_i,\T_i)\to (X_i/\cc{R}_i,\T/\cc{R}_i)$ es abierta. Demostrar que
            \begin{equation*}
                (X_1\times X_2)/\cc{R} \cong (X_1/\cc{R}_1,\T_1/\cc{R}_1)\times (X_2/\cc{R}_2,\T_2/\cc{R}_2),
            \end{equation*}
            donde $\cc{R}$ es la relación de equivalencia en $X_1\times X_2$ definida por
            \begin{equation*}
                (x_1,x_2)\cc{R}(y_1,y_2) \iff x_1\cc{R}_1 y_1 \text{ y } x_2\cc{R}_2 y_2.
            \end{equation*}

            Sea $f=p_1\times p_2: (X_1\times X_2, \T_1\times \T_2)\to (X_1/\cc{R}_1,\T_1/\cc{R}_1)\times (X_2/\cc{R}_2,\T_2/\cc{R}_2)$
            aplicación producto dada por $f(x,y)=(p_1(x), p_2(y))$. Como $p_1, p_2$ son abiertas, continuas y sobreyectivas, tenemos que $f$ es continua, abierta y sobreyectiva,
            por lo que $f$ es una identificación. Veamos que $\cc{R}=\cc{R}_f$:
            \begin{align*}
                (x,y)\cc{R}_f (x',y') &\iff
                f(x,y)=f(x',y') \iff
                (p_1(x), p_2(y)) = (p_1(x'), p_2(y')) \iff\\
                &\iff p_1(x)=p_1(x') \text{ y } p_2(y)=p_2(y') \iff
                x\cc{R}_1 x' \text{ y } y\cc{R}_2 y' \iff\\
                &\iff (x,y)\cc{R} (x',y')
            \end{align*}

            Por tanto, como $f$ es una identificación y $\cc{R}=\cc{R}_f$, tenemos que $f$ induce un homeomorfismo entre $(X_1\times X_2)/\cc{R}$ y $(X_1/\cc{R}_1,\T_1/\cc{R}_1)\times (X_2/\cc{R}_2,\T_2/\cc{R}_2)$.


            \item Sea $f:(X,\T)\to (Y,\T')$ una aplicación entre espacios topológicos tal que $f(A)$ es conexo en $(Y,\T')$
            para cada $A$ conexo en $(X,\T)$. ¿Es $f$ continua?\\

            Sea $X=\{1,2\}$, y consideramos las topologías de Sierpinsky,
            $\T_1 = \{\emptyset, X, \{1\}\}$ y $\T_2 = \{\emptyset, X, \{2\}\}$. Sea ahora ${Id}_X:\left(X, \T_1\right)\to \left(X, \T_2\right)$.
            Claramente $\{1\},\{2\}$ son conexos en ambas topologías porque no hay dos abiertos no triviales,
            por lo que $f(A)$ es conexo en $(X,\T_2)$ para cada $A$ conexo en $(X,\T_1)$. Sin embargo, $f$ no es continua, ya que
            $f^{-1}(\{2\}) = \{2\}\notin \T_1$.
        \end{enumerate}
    \end{ejercicio}
\end{document}